Incumbency protection is one of the most commonly cited justifications for departures
from neutral redistricting criteria, yet its legal status in federal and state law
is poorly understood by practitioners on both sides of redistricting disputes. This
paper provides a comprehensive legal analysis of incumbency protection as a redistricting
criterion: its permissibility under federal constitutional doctrine, its interaction
with the Voting Rights Act of 1965, and the state constitutional frameworks that govern
its use. We identify three federal cases that define the constitutional parameters:
\textit{Karcher v.\ Daggett} (1983), which identifies incumbency protection as a
legitimate objective that may justify minor population deviations; \textit{Cox v.\
Larios} (2004), which limits that permission to deviations not exceeding constitutional
thresholds; and \textit{Thornburg v.\ Gingles} (1986), which creates a tension between
incumbency protection and VRA \S{}2 compliance in majority-minority districts. We then
analyse the implications of this doctrine for the \textsc{bisect} algorithmic redistricting
system, which neither protects nor targets incumbents. Our central legal conclusion is
that \textsc{bisect}'s constitutional silence on incumbency---the algorithm receives no
incumbent address data and produces incumbency outcomes purely as geometric consequences
of population and compactness optimisation---is legally sound under all three levels of
the analysis: federal constitutional doctrine does not require incumbent protection;
VRA compliance does not require incumbent addresses; and state constitutional frameworks
that permit incumbency consideration do not mandate it. The algorithmic map's incumbency
outcomes (documented in I.1--I.3) are indistinguishable from what a neutral commission
instructed to ignore incumbency would produce, establishing \textsc{bisect} as a
constitutionally defensible redistricting methodology on incumbency grounds.
